% !TEX program = xelatex
% TP1 CampusMarket - Rapport LaTeX
\documentclass[12pt,a4paper,oneside]{report}

% === INFORMATIONS TP ===
\newcommand{\ecoleNom}{École Nationale Supérieure de l'Intelligence Artificielle et Sciences des Données}
\newcommand{\ecoleVille}{ENSIASD Taroudant}
\newcommand{\filiere}{Génie logiciel}

\newcommand{\moduleNom}{Gestion de Projet \& Génie Logiciel}
\newcommand{\tpNumero}{TP N\textsuperscript{o}~1}
\newcommand{\tpTitre}{CampusMarket -- Plateforme d'achat et de vente de produits de seconde main entre étudiants}

\newcommand{\etudiantNom}{Équipe CampusMarket}
\newcommand{\etudiantGroupe}{Groupe B}
\newcommand{\enseignantNom}{Pr. Latifa RASSAM}
\newcommand{\dateRemise}{16 octobre 2026}
\newcommand{\anneeUniv}{2025 / 2026}

\newcommand{\logoEcole}{images/logo_ensiasd}

% === PREAMBULE ===
\usepackage[a4paper,top=2.5cm,bottom=2.5cm,left=2.5cm,right=2.5cm]{geometry}
\usepackage{amsmath,amssymb}
\usepackage{xcolor}

\definecolor{couvbleu}{RGB}{41,84,150}

\usepackage{graphicx}
\usepackage{tikz}
\usepackage{caption}
\usepackage{titlesec}
\usepackage{fancyhdr}
\usepackage{ragged2e}
\usepackage{enumitem}
\usepackage{etoolbox}
\usepackage{url}
\usepackage{setspace}
\onehalfspacing

\setlength{\emergencystretch}{2em}

\usepackage{fontspec}
\setmainfont{TeX Gyre Termes}
\setmonofont{TeX Gyre Cursor}[Scale=0.9]

% === PACKAGES PERSONNELS ===
\usepackage{pgfgantt}
\usepackage{booktabs}
\usepackage{array}
\usepackage{longtable}
\usepackage{tabularx}
\usepackage{pdflscape}
\usepackage{ragged2e}


% === VARIABLES PROJET ===
% === PROJET ===
\newcommand{\nomProjet}{CampusMarket}
\newcommand{\intitule}{Plateforme d\'achat et de vente de produits de seconde main entre \'etudiants}
\newcommand{\module}{Gestion de Projet \& G\'enie Logiciel}
\newcommand{\tpNom}{TP1 : Lancement, analyse et planification d’un projet logiciel}

% === EQUIPE ===
\newcommand{\equipeMembres}{%
  Membre 1&Product Owner\\
  Membre 2&Scrum Master\\
  Membre 3&Tech Lead\\
  Membre 4&D\'eveloppeur\\
  Membre 5&D\'eveloppeur\\
  Membre 6&D\'eveloppeur%
}

% === SEANCES ===
\newcommand{\seancesLib}{S1 -- S6}

% === CHEMIN CRITIQUE ===
\newcommand{\cheminCritique}{T1 $\to$ T2 $\to$ T3 $\to$ T4 $\to$ T5 $\to$ T6 $\to$ T7 $\to$ T10 $\to$ T11}

% === CONTACT MVP ===
\newcommand{\methodeContactMVP}{Affichage des coordonn\'ees du vendeur (t\'el\'ephone / WhatsApp / e-mail) renseign\'ees \`a la publication}


% === CONFIG ===
% Configurations personnelles (facultatives pour TP1)


% === LANGUE FRANCAISE ===
\usepackage{polyglossia}
\setdefaultlanguage{french}

% === STYLES TITRES ===
\titleformat{\chapter}[block]
  {\normalfont\bfseries}
  {\fontsize{16}{20}\selectfont\MakeUppercase{Chapitre \thechapter}~:}
  {0.4em}
  {\raggedright\normalfont\fontsize{16}{20}\selectfont\bfseries\MakeUppercase}
\titlespacing*{\chapter}{0pt}{0pt}{24pt}

\renewcommand{\thechapter}{\arabic{chapter}}

\titleformat{\section}
  {\normalfont\fontsize{14}{17}\selectfont\bfseries}
  {\thesection}
  {0.6em}
  {}
\titlespacing*{\section}{0pt}{12pt}{10pt}

\titleformat{\subsection}
  {\normalfont\fontsize{12}{15}\selectfont\bfseries}
  {\thesubsection}
  {0.6em}
  {}
\titlespacing*{\subsection}{0pt}{8pt}{6pt}

\renewcommand{\thesection}{\thechapter.\arabic{section}}
\renewcommand{\thesubsection}{\thesection.\arabic{subsection}}

\setcounter{secnumdepth}{2}

% === TABLEAUX/FIGURES ===
\renewcommand{\thefigure}{\thechapter.\arabic{figure}}
\renewcommand{\thetable}{\thechapter.\arabic{table}}

\captionsetup[figure]{font={small,it}, justification=centering, labelfont=bf, labelsep=space, skip=6pt}
\captionsetup[table]{font=small, justification=centering, labelfont=bf, labelsep=space, skip=6pt}

% === EN-TETES/PAGES ===
\pagestyle{fancy}
\fancyhf{}
\fancyfoot[R]{\thepage}
\renewcommand{\headrulewidth}{0pt}
\fancypagestyle{plain}{\fancyhf{}\fancyfoot[R]{\thepage}\renewcommand{\headrulewidth}{0pt}}

% === HYPERREF ===
\usepackage[colorlinks=true, linkcolor=couvbleu, citecolor=couvbleu, urlcolor=couvbleu]{hyperref}

% === BANDE BLEUE ===
\newcommand{\bandebleue}[1]{%
  \par\noindent
  \hspace*{-\dimexpr\oddsidemargin+1in\relax}%
  \textcolor{couvbleu}{\rule{\paperwidth}{#1}}%
  \par
}

% === DOCUMENT ===
\begin{document}

% PAGE DE GARDE
\pagestyle{empty}
\thispagestyle{empty}
\vspace*{-1cm}
\begin{center}
  \IfFileExists{images/logo_ensiasd.png}{\includegraphics[width=0.5\textwidth]{images/logo_ensiasd.png}}{\textcolor{gray}{[logo ENSIASD]}}
\end{center}
\vspace{0.2cm}
\begin{center}
  {\small\ecoleNom}
\end{center}
\vspace{0.3cm}
\bandebleue{0.32cm}
\vspace{0.6cm}
\begin{spacing}{1.2}
\begin{center}
  {\fontsize{18}{22}\selectfont\bfseries\MakeUppercase{Compte Rendu de Travaux Pratiques}}
  \vspace{0.4cm}
  {\bfseries Module~: «~\moduleNom~»}
  \vspace{0.6cm}
  \textcolor{couvbleu}{\rule{\linewidth}{1mm}}
  \vspace{0.3cm}
  {\Large\bfseries \tpNumero}
  \vspace{0.2cm}
  {\fontsize{16}{20}\selectfont\bfseries \tpTitre}
  \vspace{0.3cm}
  \textcolor{couvbleu}{\rule{\linewidth}{1mm}}
\end{center}
\end{spacing}
\vspace{0.8cm}
\begin{spacing}{1.15}
\noindent
\begin{tabularx}{\linewidth}{@{}>{\bfseries}p{4.5cm}X@{}}
  Filière~:        & \filiere \\[4pt]
  Réalisé par~:    & \etudiantNom \\[4pt]
  Groupe~:         & \etudiantGroupe \\[4pt]
  Encadré par~:    & \enseignantNom \\[4pt]
  Date de remise~: & \dateRemise \\
\end{tabularx}
\end{spacing}
\vfill
\begin{center}
  \textbf{Année Universitaire~: \anneeUniv}
\end{center}
\vspace{0.3cm}
\bandebleue{0.32cm}
\clearpage

% SOMMAIRE
\pagestyle{plain}
\pagenumbering{Roman}
\phantomsection
\addcontentsline{toc}{chapter}{Sommaire}
\tableofcontents
\clearpage

\phantomsection
\addcontentsline{toc}{chapter}{Liste des tableaux}
\listoftables
\clearpage

% CORPS
\pagenumbering{arabic}
\setcounter{page}{1}
\pagestyle{fancy}

\chapter{Présentation du projet}

Dans le cadre du module \module, nous nous intéressons à la plateforme \nomProjet\ (\intitule).

Les étudiants ont souvent besoin de livres, de calculatrices, de fournitures scolaires et d'autres équipements pour leurs études. Cependant, l'achat de produits neufs peut représenter une dépense importante pour certains d'entre eux.

Parallèlement, d'autres étudiants possèdent des produits encore utilisables dont ils n'ont plus besoin et qu'ils souhaitent vendre.

Actuellement, les annonces sont dispersées entre les groupes WhatsApp, les réseaux sociaux et le bouche-à-oreille, ce qui rend difficile la recherche de produits, la comparaison des offres et la mise en relation entre acheteurs et vendeurs.

L'objectif est de centraliser ces annonces dans un espace unique afin de faciliter l'achat et la vente de produits de seconde main entre étudiants. Cette idée devra être confirmée auprès de quelques étudiants afin de vérifier qu'elle correspond à une difficulté réellement rencontrée.

\chapter{Problématique}

La dispersion des annonces entre différents canaux (groupes WhatsApp, réseaux sociaux, bouche-à-oreille) complique la recherche de produits de seconde main pour les étudiants.

Par conséquent, la problématique du projet est la suivante :

\begin{quote}
\textit{Comment faciliter l'achat et la vente de produits de seconde main entre étudiants au sein d'une plateforme centralisée, accessible et adaptée à leurs besoins ?}
\end{quote}

Cette problématique devra être confirmée auprès de quelques étudiants afin de vérifier qu'elle correspond à une difficulté réellement rencontrée.

\chapter{Objectifs du projet}

\section{Objectif général}

L'objectif général du projet est de développer \nomProjet, une plateforme web centralisée permettant aux étudiants d'acheter et de vendre des produits de seconde main.

\section{Objectifs spécifiques}

Les objectifs spécifiques sont :

\begin{itemize}
  \item Faciliter la vente et l'achat de produits de seconde main entre étudiants.
  \item Permettre aux étudiants de trouver des produits à des prix plus accessibles.
  \item Offrir aux étudiants un espace centralisé pour publier et consulter des annonces.
  \item Faciliter la mise en relation entre acheteurs et vendeurs.
  \item Encourager la réutilisation des produits encore fonctionnels.
  \item Développer une application fonctionnelle, simple à utiliser et réalisable dans le temps disponible.
\end{itemize}

\chapter{Parties prenantes et utilisateurs ciblés}

\section{Parties prenantes}

Le tableau ci-dessous présente les principales parties prenantes du projet.

\begin{table}[h]
\centering
\begin{tabularx}{\textwidth}{>{\bfseries}l X}
\toprule
\textbf{Partie prenante} & \textbf{Rôle / Intérêt} \\
\midrule
Étudiant acheteur & Recherche des produits, consulte les annonces et souhaite acquérir des articles à moindre coût. \\
Étudiant vendeur & Publie des annonces, modifie ses produits et souhaite valoriser des objets dont il n'a plus besoin. \\
Administrateur & Gère les utilisateurs et modère les annonces signalées. \\
Équipe de développement & Analyse, conçoit, développe et teste la plateforme. \\
Enseignante & Encadre et évalue le projet dans le cadre du module. \\
Communauté étudiante & Bénéficie de la réutilisation des produits et d'un marché de seconde main plus accessible. \\
\bottomrule
\end{tabularx}
\caption{Parties prenantes du projet}
\end{table}

\section{Acteurs du système}

Les acteurs sont les utilisateurs qui interagissent directement avec le système.

\begin{table}[h]
\centering
\begin{tabularx}{\textwidth}{>{\bfseries}l X}
\toprule
\textbf{Acteur} & \textbf{Description} \\
\midrule
Étudiant acheteur & Recherche des produits, consulte les annonces et contacte les vendeurs. \\
Étudiant vendeur & Publie des annonces, modifie ses propres annonces et gère leur disponibilité. \\
Administrateur & Gère les utilisateurs et modère les annonces signalées. \\
\bottomrule
\end{tabularx}
\caption{Acteurs du système}
\end{table}

\textit{Remarque :} Une partie prenante est une personne ou une organisation concernée par le projet, tandis qu'un acteur interagit directement avec le système.

\chapter{Besoins fonctionnels et non fonctionnels identifiés}

\section{Besoins fonctionnels}

Le tableau ci-dessous présente les besoins fonctionnels (BF) du système.

\begin{table}[ht]
\centering
\renewcommand{\arraystretch}{1.25}
\begin{tabularx}{\textwidth}{>{\bfseries}c X}
\toprule
\textbf{ID} & \textbf{Besoin fonctionnel} \\
\midrule
BF-01 & Le système doit permettre à un étudiant de créer un compte et de se connecter. \\
BF-02 & Le système doit permettre à un étudiant de publier une annonce contenant un titre, une description, un prix, une catégorie et des photos. \\
BF-03 & Le système doit permettre à un vendeur de gérer ses propres annonces (consulter, modifier ou supprimer). \\
BF-04 & Le système doit permettre à un acheteur de consulter les annonces disponibles et de consulter les détails d'un produit (fiche produit) avec les informations fournies par le vendeur. \\
BF-05 & Le système doit permettre à un acheteur de rechercher et de filtrer les annonces (par mots-clés, par catégorie et par fourchette de prix). \\
BF-06 & Le système doit permettre à un acheteur d'entrer en contact externe avec le vendeur (via les coordonnées fournies : téléphone, WhatsApp ou e-mail) pour obtenir des informations sur un produit. \\
BF-07 & Le système doit permettre à un vendeur de marquer une annonce comme réservée ou vendue. \\
BF-08 & Le système doit permettre à un utilisateur de signaler une annonce inappropriée ou suspecte. \\
BF-09 & Le système doit permettre à un administrateur de consulter les annonces signalées, de les modérer, de supprimer une annonce inappropriée et de bannir un utilisateur en cas d'abus. \\
\bottomrule
\end{tabularx}
\caption{Besoins fonctionnels}
\end{table}

\section{Besoins non fonctionnels}

Le tableau ci-dessous présente les besoins non fonctionnels (BNF) du système.

\begin{table}[ht]
\centering
\renewcommand{\arraystretch}{1.25}
\begin{tabularx}{\textwidth}{>{\bfseries}c c X}
\toprule
\textbf{ID} & \textbf{Catégorie} & \textbf{Besoin non fonctionnel} \\
\midrule
BNF-01 & Sécurité & Le système doit protéger les mots de passe et limiter l'accès aux données selon les autorisations de chaque utilisateur. \\
BNF-02 & Performance & Les pages principales et les résultats de recherche devraient s'afficher en moins de deux secondes dans les conditions de charge définies pour le projet. \\
BNF-03 & Ergonomie & L'interface doit être simple à comprendre et utilisable sur ordinateur et smartphone. \\
BNF-04 & Maintenabilité & Le code doit être organisé en modules distincts afin de faciliter les corrections et les évolutions. \\
\bottomrule
\end{tabularx}
\caption{Besoins non fonctionnels}
\end{table}

\chapter{Périmètre du projet}

\section{Périmètre inclus}

Le projet se concentre sur un MVP (Minimum Viable Product) permettant d'assurer les fonctionnalités essentielles pour l'achat et la vente de produits de seconde main entre étudiants.

Les fonctionnalités prioritaires incluses dans le MVP sont :

\begin{table}[h]
\centering
\renewcommand{\arraystretch}{1.25}
\begin{tabularx}{\textwidth}{>{\bfseries}c X c}
\toprule
\textbf{ID} & \textbf{Fonctionnalité} & \textbf{Priorité} \\
\midrule
F1 & Inscription et connexion des utilisateurs & Haute \\
F2 & Publication et gestion des annonces & Haute \\
F3 & Consultation du catalogue des produits & Haute \\
F4 & Recherche et filtrage des annonces & Haute \\
F5 & Consultation des détails d'un produit & Haute \\
F6 & Prise de contact entre acheteur et vendeur & Haute \\
F7 & Gestion de la disponibilité des produits (réservé / vendu) & Haute \\
F8 & Signalement et modération des annonces & Moyenne \\
F9 & Gestion du profil utilisateur & Moyenne \\
F10 & Consultation des annonces publiées par un utilisateur & Moyenne \\
\bottomrule
\end{tabularx}
\caption{Fonctionnalités incluses dans le MVP}
\end{table}

\section{Périmètre exclu}

Afin de rester réaliste compte tenu du temps disponible (six séances de TP), certaines fonctionnalités ne seront pas implémentées dans cette première version.

\begin{table}[h]
\centering
\renewcommand{\arraystretch}{1.25}
\begin{tabularx}{\textwidth}{>{\bfseries}X X}
\toprule
\textbf{Fonctionnalité exclue} & \textbf{Justification} \\
\midrule
Paiement en ligne & Nécessite une intégration de paiement et des mesures de sécurité supplémentaires. \\
Livraison intégrée et suivi des colis & Implique une logistique et des services externes qui dépassent le périmètre du prototype. \\
Application mobile native Android/iOS & Une application web responsive suffit pour la première version. \\
Système avancé de notation et de réputation & Peut être ajouté dans une version ultérieure après validation du fonctionnement principal. \\
Recommandations personnalisées par intelligence artificielle & Cette fonctionnalité n'est pas essentielle au besoin principal et augmenterait la complexité. \\
\bottomrule
\end{tabularx}
\caption{Fonctionnalités exclues du périmètre}
\end{table}

\section{Justification du périmètre}

La définition du périmètre permet de limiter la complexité, de mieux estimer les tâches et de concentrer les efforts de l'équipe sur les fonctionnalités prioritaires dans le temps imparti. \nomProjet\ vise ainsi à répondre au besoin réel des étudiants tout en restant réalisable dans le cadre de ce projet.

\chapter{Backlog initial}

Le backlog initial est présenté ci-dessous selon les phases retenues pour le projet.

\begin{table}[ht]
\centering
\renewcommand{\arraystretch}{1.25}
\small
\begin{tabularx}{\textwidth}{>{\bfseries}c X c c X}
\toprule
\textbf{ID} & \textbf{Tâche} & \textbf{Semaine} & \textbf{Prédécesseurs} & \textbf{Responsable principal} \\
\midrule
T1 & Analyse : identification et validation des besoins, analyse des acteurs, rédaction des exigences et définition du périmètre. & Semaine 1 & -- & Product Owner \\
T2 & Conception : préparation des éléments de conception. & Semaine 2 & T1 & Architecte \\
T3 & Développement backend : mise en place de l'environnement, authentification, gestion des utilisateurs, annonces, catalogue, recherche et filtres, disponibilité. & Semaine 3 & T2 & Développeur Backend \\
T4 & Développement backend : suite du développement backend et finalisation de la logique métier. & Semaine 4 & T3 & Développeur Backend \\
T5 & Développement frontend : interfaces utilisateur, intégration avec le backend, prise de contact externe et modération. & Semaine 5 & T4 & Développeur Frontend / Équipe \\
T6 & Tests et Validation : tests des fonctionnalités, corrections et validation. & Semaine 6 & T5 & Toute l'équipe \\
T7 & Déploiement et Maintenance : déploiement de l'application et actions de maintenance préliminaires. & Semaine 6 & T6 & Tech Lead / Équipe \\
T8 & Présentation : préparation et présentation du projet. & Semaine 7 & T7 & Scrum Master / Toute l'équipe \\
\bottomrule
\end{tabularx}
\caption{Backlog initial}
\end{table}

\chapter{Planification prévisionnelle du projet (Gantt) et chemin critique}

\section{Planification par phases}

La planification s'étend sur six semaines de réalisation, avec une septième semaine consacrée exclusivement à la préparation et à la présentation du projet.

Les phases retenues sont les suivantes :

\begin{enumerate}
  \item \textbf{Analyse} (T1) : identification et validation des besoins, analyse des acteurs, rédaction des exigences et définition du périmètre.
  \item \textbf{Conception} (T2) : préparation des éléments de conception.
  \item \textbf{Développement backend} (T3 -- T4) : développement de la partie serveur (authentification, gestion des utilisateurs, annonces, catalogue, recherche/filtres, disponibilité).
  \item \textbf{Développement frontend} (T5) : développement des interfaces et intégration avec le backend (contact externe, modération).
  \item \textbf{Tests et Validation} (T6) : tests des fonctionnalités, corrections et validation.
  \item \textbf{Déploiement et Maintenance} (T7) : déploiement et actions de maintenance préliminaires.
  \item \textbf{Présentation} (T8) : préparation et présentation du projet (Semaine 7).
\end{enumerate}

\section{Tâches, durées, prédécesseurs et responsables}

Le tableau ci-dessous présente le détail des tâches.

\begin{table}[ht]
\centering
\renewcommand{\arraystretch}{1.25}
\small
\begin{tabularx}{\textwidth}{>{\bfseries}c X c c X}
\toprule
\textbf{ID} & \textbf{Tâche} & \textbf{Durée estimée} & \textbf{Prédécesseurs} & \textbf{Responsable principal} \\
\midrule
T1 & Analyse : besoins, acteurs, exigences et périmètre & 1 semaine & -- & Product Owner \\
T2 & Conception & 1 semaine & T1 & Architecte \\
T3 & Développement backend -- Partie 1 & 1 semaine & T2 & Développeur Backend \\
T4 & Développement backend -- Partie 2 & 1 semaine & T3 & Développeur Backend \\
T5 & Développement frontend & 1 semaine & T4 & Développeur Frontend / Équipe \\
T6 & Tests et Validation & 1 semaine & T5 & Toute l'équipe \\
T7 & Déploiement et Maintenance & 1 semaine & T6 & Tech Lead / Équipe \\
T8 & Présentation : préparation et présentation du projet & 1 semaine (S7) & T7 & Scrum Master / Toute l'équipe \\
\bottomrule
\end{tabularx}
\caption{Tâches, durées, prédécesseurs et responsables}
\end{table}

\section{Planification prévisionnelle (Diagramme de Gantt)}

La figure ci-dessous présente la planification prévisionnelle sur 7 semaines (S1 -- S7).

\begin{ganttchart}[
  hgrid,
  vgrid,
  x unit=1.6cm,
  y unit chart=0.55cm,
  bar/.append style={fill=blue!30},
  bar incomplete/.append style={fill=blue!15},
  title/.append style={fill=gray!10},
  title label font=\small\bfseries,
  bar label font=\small,
  group label font=\small\bfseries
]{1}{7}
  \gantttitle{S1}{1}
  \gantttitle{S2}{1}
  \gantttitle{S3}{1}
  \gantttitle{S4}{1}
  \gantttitle{S5}{1}
  \gantttitle{S6}{1}
  \gantttitle{S7}{1} \\

  \ganttbar{T1 - Analyse}{1}{1} \\
  \ganttbar{T2 - Conception}{2}{2} \\
  \ganttbar{T3 - Dév. backend Partie 1}{3}{3} \\
  \ganttbar{T4 - Dév. backend Partie 2}{4}{4} \\
  \ganttbar{T5 - Dév. frontend}{5}{5} \\
  \ganttbar{T6 - Tests et Validation}{6}{6} \\
  \ganttbar{T7 - Déploiement et Maintenance}{6}{6} \\
  \ganttbar{T8 - Présentation}{7}{7}
\end{ganttchart}
\caption{Planning prévisionnel (Gantt) sur 7 semaines (S1--S7)}
\label{gantt:planning}

\section{Chemin critique}

Le chemin critique déterminant la durée minimale du projet est le suivant :

\begin{center}
\textbf{T1 $\to$ T2 $\to$ T3 $\to$ T4 $\to$ T5 $\to$ T6 $\to$ T7 $\to$ T8}
\end{center}

\textbf{Interprétation :} tout retard sur une tâche appartenant à ce chemin critique entraînerait, sans réorganisation du travail, un retard équivalent sur la livraison finale.

\chapter{Principaux risques identifiés}

\begin{table}[ht]
\centering
\renewcommand{\arraystretch}{1.25}
\small
\begin{tabularx}{\textwidth}{>{\bfseries}c X c c c X}
\toprule
\textbf{ID} & \textbf{Risque} & \textbf{P} & \textbf{I} & \textbf{P×I} & \textbf{Mesure préventive} \\
\midrule
R1 & Manque de temps pour développer toutes les fonctionnalités & 4 & 5 & 20 & Prioriser le MVP et respecter le découpage par phases. \\
R2 & Mauvaise compréhension des besoins réels des étudiants & 3 & 4 & 12 & Interroger quelques étudiants avant de finaliser les exigences et le périmètre. \\
R3 & Difficultés techniques lors de l'intégration & 3 & 4 & 12 & Intégrer régulièrement et bien définir les interfaces entre backend et frontend. \\
R4 & Indisponibilité d'un membre de l'équipe & 3 & 3 & 9 & Partager les connaissances et documenter le travail. \\
R5 & Retard dans la conception de la base de données & 3 & 4 & 12 & Valider tôt le modèle de données et les relations principales. \\
R6 & Annonces fausses ou contenu inapproprié & 3 & 2 & 6 & Mettre en œuvre le signalement, la modération, la suppression d'annonce et le bannissement d'utilisateur. \\
R7 & Dérive du périmètre (scope creep) & 3 & 4 & 12 & Se tenir strictement dans le périmètre du MVP. \\
R8 & Indécision sur le choix technologique (framework / base de données) & 3 & 4 & 12 & Trancher sur les choix technologiques dès la phase de conception. \\
\bottomrule
\end{tabularx}
\caption{Principaux risques identifiés}
\end{table}

\section{Matrice des risques}

\textit{Matrice des risques à compléter ultérieurement.}

\chapter{Conclusion}

\addcontentsline{toc}{chapter}{Conclusion}

\textbf{\nomProjet} vise à répondre à un besoin potentiel de la communauté étudiante : faciliter l'achat et la vente de produits de seconde main dans un espace centralisé, accessible et adapté à leurs attentes.

Le projet propose un périmètre réaliste pour cette première version (MVP), comprenant la gestion des utilisateurs, la publication et la consultation des annonces, la recherche et le filtrage, la prise de contact entre acheteurs et vendeurs, ainsi qu'un mécanisme minimal de signalement et de modération. Les paiements en ligne, la livraison intégrée, l'échange sous forme de troc et les fonctionnalités avancées sont, quant à eux, reportés à des versions futures.

La réussite du projet dépendra principalement de la validation des besoins auprès des étudiants, d'une bonne priorisation des fonctionnalités et d'une planification rigoureuse sur les six séances de TP. En se concentrant sur l'essentiel et en restant vigilant vis-à-vis du chemin critique, l'équipe sera en mesure de livrer une première version fonctionnelle et cohérente avec les objectifs fixés.


\end{document}
